\documentclass[11pt,paper=a4,answers]{exam}
\usepackage{graphicx,lastpage}
\usepackage{upgreek}
\usepackage{censor}
\usepackage{gensymb}
\usepackage{amsmath}
\usepackage{multicol}
\usepackage{enumitem}
\usepackage{tasks}
\usepackage{adjustbox}
\usepackage{xcolor}
\usepackage{tfrupee}
\setlength{\voffset}{0.25in}
\usepackage{tabularx}
\newcolumntype{Y}{>{\centering\arraybackslash}X}
%==============================================================
% Customise Non-Essential Packages Below
% =============================================================
\usepackage[most]{tcolorbox}
\usepackage{eso-pic}
\usepackage{tikz}
\AddToShipoutPictureBG{%
\begin{tikzpicture}[remember picture,overlay]
\foreach \x in {-8,-4,0,4,8,12,16,20}
\foreach \y in {-8,-4,0,4,8,12,16,20} {
\node[opacity=0.05,rotate=45,anchor=center] at ([xshift=\x cm,yshift=\y cm]current page.center) {%
\begin{minipage}{2cm}
\centering
\includegraphics[width=2cm]{images/logo_black.png}\\
\small preppeo.com
\end{minipage}%
};
}
\end{tikzpicture}%
}
\printanswers

%==============================================================
% Customise Non-Essential Packages Above
% =============================================================
\definecolor{svgray}{rgb}{0.90, 0.90, 0.90}
\graphicspath{{images/}}

\usepackage[normalem]{ulem}
\renewcommand\ULthickness{1pt}   %%---> For changing thickness of underline
\setlength\ULdepth{3.5ex}%\maxdimen ---> For changing depth of underline
\renewcommand{\baselinestretch}{1}
\chead{
\includegraphics[scale=0.1,height=2cm ]{logo.png}
}
\pagestyle{empty}

\pagestyle{headandfoot}
\headrule
\cfoot{W: https://preppeo.com; M: +91-8130711689}

\begin{document}
%==============================================================
% Customise Question paper details below this
% =============================================================
\begin{center}
    \centering
    \renewcommand{\arraystretch}{1.5}
    \begin{tabularx}{\textwidth}{YYY}
    & \normalsize\bf{{\Large{{Quantitative Reasoning}}}} & \\
    By Preppeo & \normalsize\bf{{sat\_lid\_015}} & SAT\\
    \normalsize{{\textbf{{Unit:}} Advanced Math}} & \normalsize{{\textbf{{Chapter:}} Nonlinear equations in one variable}} & \normalsize{{\textbf{{Topic:}} Discriminant \& Number of Solutions}} \\
    \end{tabularx}
\end{center}
\par\noindent\rule{\textwidth}{1pt}
% =============================================================
% Customise Question paper details above this
% =============================================================

\begin{questions}
\pointsinrightmargin
\pointsdroppedatright
\marksnotpoints
\pointpoints{mark}{marks}
\pointformat{\boldmath\themarginpoints}
\bracketedpoints

% =============================================================
% Question Begin
% =============================================================
% Lid Topic: Advanced Math — Linear-Quadratic Systems
% Total Questions: 50

% --- PART 1: Screenshot-Inspired Systems (1-25) ---

% Question 1
% Category: Practice | Difficulty: Medium | Question Type: Multiple Choice
\question
A system of equations consists of a quadratic function $y = x^2 - 4x + 7$ and a linear function $y = 3$. At how many points do the graphs of these two functions intersect?
\begin{tasks}(4)
    \task Zero
    \task One
    \task Two
    \task Infinitely many
\end{tasks}

\begin{solution}
\textbf{Conceptual Explanation:} \\
The points of intersection are the values of $x$ where both equations are equal. We can find this by substituting the linear equation into the quadratic equation and analyzing the resulting quadratic.

\vspace{20pt}

\textbf{Calculation and Logic:} \\
1. Set the equations equal to each other: \\
   $x^2 - 4x + 7 = 3$

\vspace{10pt}

2. Move all terms to one side to set the equation to zero: \\
   $x^2 - 4x + 4 = 0$

\vspace{10pt}

3. Analyze the discriminant ($b^2 - 4ac$): \\
   $a = 1, b = -4, c = 4$ \\
   $D = (-4)^2 - 4(1)(4) = 16 - 16 = 0$

\vspace{10pt}

4. A discriminant of 0 means there is exactly \textbf{one} real solution, which corresponds to one point of intersection.

\vspace{25pt}

Answer: \colorbox{yellow!100}{B}
\end{solution}

\vspace{40pt}

% Question 2
% Category: Practice | Difficulty: Hard | Question Type: Multiple Choice
\question
If $(x, y)$ is a solution to the system of equations below and $x > 0$, what is the value of $x$?
\[ \begin{cases} y = x^2 - 5 \\ y = 4x \end{cases} \]
\begin{tasks}(4)
    \task 1
    \task 5
    \task 4
    \task 2
\end{tasks}

\begin{solution}
\textbf{Conceptual Explanation:} \\
To find the solution $(x, y)$, use substitution to create a single quadratic equation in terms of $x$. Then, solve the quadratic and apply the constraint $x > 0$.

\vspace{20pt}

\textbf{Calculation and Logic:} \\
1. Substitute $4x$ for $y$ in the first equation: \\
   $4x = x^2 - 5$

\vspace{10pt}

2. Rearrange into standard form: \\
   $x^2 - 4x - 5 = 0$

\vspace{10pt}

3. Factor the quadratic: \\
   $(x - 5)(x + 1) = 0$

\vspace{10pt}

4. Solve for $x$: \\
   $x = 5$ or $x = -1$

\vspace{10pt}

5. Since the problem requires $x > 0$, the only valid solution is 5.

\vspace{25pt}

Answer: \colorbox{yellow!100}{B}
\end{solution}

\vspace{40pt}

% Question 3
% Category: Practice | Difficulty: Medium | Question Type: Multiple Choice
\question

The graph of the system of equations $y = x^2 - 6$ and $y = 3$ is shown in the $xy$-plane. Which of the following are the $x$-coordinates of the points of intersection?
\begin{tasks}(2)
    \task $x = 3$ and $x = -3$
    \task $x = 6$ and $x = -6$
    \task $x = 0$ and $x = 3$
    \task $x = 9$ and $x = -9$
\end{tasks}

\begin{solution}
\textbf{Conceptual Explanation:} \\
The intersection points occur where the $y$-values are the same. We set the quadratic expression equal to the constant value of the linear function.

\vspace{20pt}

\textbf{Calculation and Logic:} \\
1. Set the equations equal: \\
   $x^2 - 6 = 3$

\vspace{10pt}

2. Isolate $x^2$: \\
   $x^2 = 9$

\vspace{10pt}

3. Take the square root of both sides: \\
   $x = \pm \sqrt{9}$ \\
   $x = 3$ or $x = -3$

\vspace{25pt}

Answer: \colorbox{yellow!100}{A}
\end{solution}

\vspace{40pt}

% Question 4
% Category: Practice | Difficulty: Hard | Question Type: Multiple Choice
\question
A line with the equation $y = mx + 2$ is tangent to the parabola $y = x^2 + 6x + 11$ at exactly one point. What is one possible value for $m$?
\begin{tasks}(4)
    \task 0
    \task 2
    \task 6
    \task 12
\end{tasks}

\begin{solution}
\textbf{Conceptual Explanation:} \\
A line is tangent to a parabola if the system of equations has exactly one solution. This means the discriminant of the resulting quadratic must be zero.

\vspace{20pt}

\textbf{Calculation and Logic:} \\
1. Set the equations equal: \\
   $x^2 + 6x + 11 = mx + 2$

\vspace{10pt}

2. Rearrange into standard form ($ax^2 + bx + c = 0$): \\
   $x^2 + (6 - m)x + 9 = 0$

\vspace{10pt}

3. Identify coefficients for the discriminant: \\
   $a = 1, b = (6 - m), c = 9$

\vspace{10pt}

4. Set $D = 0$: \\
   $(6 - m)^2 - 4(1)(9) = 0$ \\
   $(6 - m)^2 - 36 = 0$ \\
   $(6 - m)^2 = 36$

\vspace{10pt}

5. Solve for $m$: \\
   $6 - m = 6 \Rightarrow m = 0$ \\
   $6 - m = -6 \Rightarrow m = 12$

\vspace{25pt}

Answer: \colorbox{yellow!100}{D}
\end{solution}

\vspace{40pt}

% Question 5
% Category: Practice | Difficulty: Medium | Question Type: Multiple Choice
\question
How many real solutions does the following system of equations have?
\[ \begin{cases} y = x^2 + 2x + 5 \\ y = -x + 1 \end{cases} \]
\begin{tasks}(4)
    \task Zero
    \task One
    \task Two
    \task Infinitely many
\end{tasks}

\begin{solution}
\textbf{Conceptual Explanation:} \\
The number of real solutions is determined by the discriminant of the quadratic formed by setting the two equations equal to each other.

\vspace{20pt}

\textbf{Calculation and Logic:} \\
1. Set the equations equal: \\
   $x^2 + 2x + 5 = -x + 1$

\vspace{10pt}

2. Move all terms to one side: \\
   $x^2 + 3x + 4 = 0$

\vspace{10pt}

3. Calculate the discriminant ($b^2 - 4ac$): \\
   $D = (3)^2 - 4(1)(4) = 9 - 16 = -7$

\vspace{10pt}

4. Since the discriminant is negative ($D < 0$), there are no real solutions.

\vspace{25pt}

Answer: \colorbox{yellow!100}{A}
\end{solution}

% Question 6
% Category: Practice | Difficulty: Medium | Question Type: Multiple Choice
\question
Which of the following points $(x, y)$ is a solution to the system of equations below?
\[ \begin{cases} y = x^2 - 2 \\ y = x \end{cases} \]
\begin{tasks}(4)
    \task $(1, 1)$
    \task $(2, 2)$
    \task $(0, 0)$
    \task $(-1, 1)$
\end{tasks}

\begin{solution}
\textbf{Conceptual Explanation:} \\
To find the solution, we set the equations equal to each other to solve for $x$. Once we find $x$, we plug it back into either equation to find the corresponding $y$-value.

\vspace{20pt}

\textbf{Calculation and Logic:} \\
1. Set the equations equal: \\
   $x^2 - 2 = x$

\vspace{10pt}

2. Rearrange to standard form: \\
   $x^2 - x - 2 = 0$

\vspace{10pt}

3. Factor the quadratic: \\
   $(x - 2)(x + 1) = 0$

\vspace{10pt}

4. Solve for $x$: \\
   $x = 2$ or $x = -1$

\vspace{10pt}

5. Find $y$ using $y = x$: \\
   If $x = 2, y = 2 \Rightarrow (2, 2)$ \\
   If $x = -1, y = -1 \Rightarrow (-1, -1)$

\vspace{25pt}

Answer: \colorbox{yellow!100}{B}
\end{solution}

\vspace{40pt}

% Question 7
% Category: Practice | Difficulty: Hard | Question Type: Multiple Choice
\question
The line $y = -2x + k$ intersects the parabola $y = x^2 - 6x + 7$ at exactly one point. What is the value of $k$?
\begin{tasks}(4)
    \task 2
    \task 3
    \task 4
    \task 7
\end{tasks}

\begin{solution}
\textbf{Conceptual Explanation:} \\
"Intersects at exactly one point" means the system has a unique solution. We set the expressions equal and ensure the discriminant of the resulting quadratic is zero.

\vspace{20pt}

\textbf{Calculation and Logic:} \\
1. Set the equations equal: \\
   $x^2 - 6x + 7 = -2x + k$

\vspace{10pt}

2. Standard form: \\
   $x^2 - 4x + (7 - k) = 0$

\vspace{10pt}

3. Identify coefficients: $a = 1, b = -4, c = (7 - k)$.

\vspace{10pt}

4. Set $D = 0$: \\
   $(-4)^2 - 4(1)(7 - k) = 0$ \\
   $16 - 4(7 - k) = 0$

\vspace{10pt}

5. Solve for $k$: \\
   $16 - 28 + 4k = 0$ \\
   $-12 + 4k = 0$ \\
   $4k = 12 \Rightarrow k = 3$

\vspace{25pt}

Answer: \colorbox{yellow!100}{B}
\end{solution}

\vspace{40pt}

% Question 8
% Category: Practice | Difficulty: Medium | Question Type: Multiple Choice
\question
If $(x, y)$ is a solution to the system of equations below, which of the following is a possible value for $y$?
\[ \begin{cases} y = x^2 \\ y = x + 6 \end{cases} \]
\begin{tasks}(4)
    \task 3
    \task 4
    \task 6
    \task 9
\end{tasks}

\begin{solution}
\textbf{Conceptual Explanation:} \\
Solve for $x$ first by substitution, then find the possible $y$-values.

\vspace{20pt}

\textbf{Calculation and Logic:} \\
1. $x^2 = x + 6 \Rightarrow x^2 - x - 6 = 0$ \\
2. $(x - 3)(x + 2) = 0 \Rightarrow x = 3$ or $x = -2$

\vspace{10pt}

3. Calculate $y$: \\
   If $x = 3, y = 3^2 = 9$ \\
   If $x = -2, y = (-2)^2 = 4$

\vspace{10pt}

4. Among the choices, both 4 and 9 are possible. (Note: SAT questions typically provide only one of the valid options). Here, we select 9.

\vspace{25pt}

Answer: \colorbox{yellow!100}{D}
\end{solution}

\vspace{40pt}

% Question 9
% Category: Practice | Difficulty: Hard | Question Type: Multiple Choice
\question
A system of equations consists of $y = (x - 3)^2 + 2$ and $y = 2$. What is the solution $(x, y)$ to the system?
\begin{tasks}(4)
    \task $(0, 2)$
    \task $(3, 2)$
    \task $(2, 3)$
    \task $(3, 0)$
\end{tasks}

\begin{solution}
\textbf{Conceptual Explanation:} \\
This system involves a parabola in vertex form and a horizontal line that passes exactly through its vertex.

\vspace{20pt}

\textbf{Calculation and Logic:} \\
1. Set $(x - 3)^2 + 2 = 2$ \\
2. $(x - 3)^2 = 0$ \\
3. $x - 3 = 0 \Rightarrow x = 3$

\vspace{10pt}

4. The $y$-value is already given as 2. \\
5. The solution is $(3, 2)$.

\vspace{25pt}

Answer: \colorbox{yellow!100}{B}
\end{solution}

\vspace{40pt}

% Question 10
% Category: Practice | Difficulty: Medium | Question Type: Multiple Choice
\question
How many points of intersection exist between the line $y = 2x - 10$ and the parabola $y = x^2 - 4x + 5$?
\begin{tasks}(4)
    \task 0
    \task 1
    \task 2
    \task 3
\end{tasks}

\begin{solution}
\textbf{Conceptual Explanation:} \\
Use the discriminant of the combined equation to check the number of real roots.

\vspace{20pt}

\textbf{Calculation and Logic:} \\
1. $x^2 - 4x + 5 = 2x - 10$ \\
2. $x^2 - 6x + 15 = 0$

\vspace{10pt}

3. $D = (-6)^2 - 4(1)(15)$ \\
4. $D = 36 - 60 = -24$

\vspace{10pt}

5. Since $D < 0$, there are zero real solutions.

\vspace{25pt}

Answer: \colorbox{yellow!100}{A}
\end{solution}

\vspace{40pt}

% Question 11
% Category: Practice | Difficulty: Hard | Question Type: Multiple Choice
\question
If $(x, y)$ is a solution to the system below and $y > 0$, what is the value of $y$?
\[ \begin{cases} y = x^2 + 1 \\ y = 2x \end{cases} \]
\begin{tasks}(4)
    \task 1
    \task 2
    \task 0
    \task No such $y$ exists.
\end{tasks}

\begin{solution}
\textbf{Conceptual Explanation:} \\
Find the intersection point and verify the $y$-coordinate.

\vspace{20pt}

\textbf{Calculation and Logic:} \\
1. $x^2 + 1 = 2x \Rightarrow x^2 - 2x + 1 = 0$ \\
2. $(x - 1)^2 = 0 \Rightarrow x = 1$

\vspace{10pt}

3. Plug $x = 1$ into $y = 2x$: \\
   $y = 2(1) = 2$

\vspace{25pt}

Answer: \colorbox{yellow!100}{B}
\end{solution}

\vspace{40pt}

% Question 12
% Category: Practice | Difficulty: Medium | Question Type: Multiple Choice
\question
The parabola $y = -x^2 + 4$ and the line $y = 0$ intersect at points $A$ and $B$. What is the distance between $A$ and $B$?
\begin{tasks}(4)
    \task 2
    \task 4
    \task 8
    \task 0
\end{tasks}

\begin{solution}
\textbf{Conceptual Explanation:} \\
The intersection points with $y = 0$ are the $x$-intercepts of the parabola. The distance between them is the difference between the larger and smaller $x$-values.

\vspace{20pt}

\textbf{Calculation and Logic:} \\
1. $-x^2 + 4 = 0 \Rightarrow x^2 = 4$ \\
2. $x = 2$ and $x = -2$

\vspace{10pt}

3. Distance $= 2 - (-2) = 4$.

\vspace{25pt}

Answer: \colorbox{yellow!100}{B}
\end{solution}

\vspace{40pt}

% Question 13
% Category: Practice | Difficulty: Hard | Question Type: Multiple Choice
\question
If $y = 3x^2 - 12x + 10$ and $y = -2$, what is the sum of the $x$-coordinates of the solutions?
\begin{tasks}(4)
    \task 4
    \task 3
    \task 12
    \task -4
\end{tasks}

\begin{solution}
\textbf{Conceptual Explanation:} \\
You can find the sum of the roots using the formula $-b/a$ after combining the equations into a standard quadratic.

\vspace{20pt}

\textbf{Calculation and Logic:} \\
1. $3x^2 - 12x + 10 = -2$ \\
2. $3x^2 - 12x + 12 = 0$

\vspace{10pt}

3. Sum $= -(-12)/3 = 4$.

\vspace{25pt}

Answer: \colorbox{yellow!100}{A}
\end{solution}

\vspace{40pt}

% Question 14
% Category: Practice | Difficulty: Medium | Question Type: Multiple Choice
\question
Which of the following describes the intersection of $y = x^2 + 10$ and $y = 5$?
\begin{tasks}(1)
    \task They intersect at two points.
    \task They intersect at one point.
    \task They do not intersect.
    \task The line is the axis of symmetry for the parabola.
\end{tasks}

\begin{solution}
\textbf{Conceptual Explanation:} \\
The vertex of the parabola $y = x^2 + 10$ is $(0, 10)$, and it opens upward. Any horizontal line below $y = 10$ will not intersect it.

\vspace{20pt}

\textbf{Calculation and Logic:} \\
1. $x^2 + 10 = 5$ \\
2. $x^2 = -5$ \\
3. Since a square cannot be negative, there are no real solutions.

\vspace{25pt}

Answer: \colorbox{yellow!100}{C}
\end{solution}

\vspace{40pt}

% Question 15
% Category: Practice | Difficulty: Hard | Question Type: Multiple Choice
\question
The system below has exactly one solution. What is the value of $c$?
\[ \begin{cases} y = x^2 + cx + 9 \\ y = 0 \end{cases} \]
\begin{tasks}(4)
    \task 3
    \task 6
    \task 9
    \task 0
\end{tasks}

\begin{solution}
\textbf{Conceptual Explanation:} \\
For the parabola to touch the $x$-axis ($y = 0$) at exactly one point, it must be a perfect square trinomial.

\vspace{20pt}

\textbf{Calculation and Logic:} \\
1. $D = c^2 - 4(1)(9) = 0$ \\
2. $c^2 - 36 = 0$ \\
3. $c^2 = 36 \Rightarrow c = 6$ (or $-6$).

\vspace{25pt}

Answer: \colorbox{yellow!100}{B}
\end{solution}

\vspace{40pt}

% Question 16
% Category: Practice | Difficulty: Medium | Question Type: Multiple Choice
\question
What is the $y$-coordinate of the intersection of $y = x^2 - x$ and $y = x + 3$?
\begin{tasks}(4)
    \task 3
    \task 6
    \task 0
    \task 1
\end{tasks}

\begin{solution}
\textbf{Conceptual Explanation:} \\
Find $x$ first, then use the linear equation to find $y$.

\vspace{20pt}

\textbf{Calculation and Logic:} \\
1. $x^2 - x = x + 3 \Rightarrow x^2 - 2x - 3 = 0$ \\
2. $(x - 3)(x + 1) = 0 \Rightarrow x = 3, -1$

\vspace{10pt}

3. If $x = 3, y = 3 + 3 = 6$. \\
4. If $x = -1, y = -1 + 3 = 2$. \\
5. Choice B is 6.

\vspace{25pt}

Answer: \colorbox{yellow!100}{B}
\end{solution}

\vspace{40pt}

% Question 17
% Category: Practice | Difficulty: Hard | Question Type: Multiple Choice
\question
At what points does the line $y = x + 1$ intersect the circle $x^2 + y^2 = 5$?
\begin{tasks}(2)
    \task $(1, 2)$ and $(-2, -1)$
    \task $(2, 1)$ and $(-1, -2)$
    \task $(0, 1)$ and $(1, 0)$
    \task $(1, 2)$ and $(2, 1)$
\end{tasks}

\begin{solution}
\textbf{Conceptual Explanation:} \\
Substitute the linear expression for $y$ into the circle's equation.

\vspace{20pt}

\textbf{Calculation and Logic:} \\
1. $x^2 + (x + 1)^2 = 5$ \\
2. $x^2 + x^2 + 2x + 1 = 5$ \\
3. $2x^2 + 2x - 4 = 0$ \\
4. $x^2 + x - 2 = 0 \Rightarrow (x + 2)(x - 1) = 0$

\vspace{10pt}

5. $x = -2, 1$. \\
6. If $x = -2, y = -1$. If $x = 1, y = 2$.

\vspace{25pt}

Answer: \colorbox{yellow!100}{A}
\end{solution}

\vspace{40pt}

% Question 18
% Category: Practice | Difficulty: Medium | Question Type: Multiple Choice
\question
If $y = x^2 - 4x + 10$ and $y = 3x$, what is the value of $x^2 - 7x$?
\begin{tasks}(4)
    \task 10
    \task -10
    \task 7
    \task 0
\end{tasks}

\begin{solution}
\textbf{Conceptual Explanation:} \\
Combine the equations and see if the resulting expression matches the target expression.

\vspace{20pt}

\textbf{Calculation and Logic:} \\
1. $x^2 - 4x + 10 = 3x$ \\
2. $x^2 - 7x + 10 = 0$ \\
3. $x^2 - 7x = -10$.

\vspace{25pt}

Answer: \colorbox{yellow!100}{B}
\end{solution}

\vspace{40pt}

% Question 19
% Category: Practice | Difficulty: Hard | Question Type: Multiple Choice
\question
The line $x = 2$ intersects the parabola $y = 3x^2 - 5x + 1$ at point $P$. What is the $y$-coordinate of $P$?
\begin{tasks}(4)
    \task 1
    \task 3
    \task 23
    \task 11
\end{tasks}

\begin{solution}
\textbf{Conceptual Explanation:} \\
A vertical line $x = c$ always intersects a function at exactly one point. Simply substitute the value of $x$ into the function.

\vspace{20pt}

\textbf{Calculation and Logic:} \\
1. $y = 3(2)^2 - 5(2) + 1$ \\
2. $y = 3(4) - 10 + 1$ \\
3. $y = 12 - 10 + 1 = 3$.

\vspace{25pt}

Answer: \colorbox{yellow!100}{B}
\end{solution}

\vspace{40pt}

% Question 20
% Category: Practice | Difficulty: Medium | Question Type: Multiple Choice
\question
For what value of $a$ will the parabola $y = ax^2$ pass through the intersection of $y = x$ and $x = 2$?
\begin{tasks}(4)
    \task 0.5
    \task 1
    \task 2
    \task 4
\end{tasks}

\begin{solution}
\textbf{Conceptual Explanation:} \\
Find the intersection point of the two lines first, then substitute that point into the parabola's equation.

\vspace{20pt}

\textbf{Calculation and Logic:} \\
1. Intersection of $y=x$ and $x=2$ is $(2, 2)$. \\
2. Substitute into $y = ax^2$: \\
   $2 = a(2^2) \Rightarrow 2 = 4a$ \\
   $a = 0.5$.

\vspace{25pt}

Answer: \colorbox{yellow!100}{A}
\end{solution}

\vspace{40pt}

% Question 21
% Category: Practice | Difficulty: Hard | Question Type: Multiple Choice
\question
The vertex of the parabola $y = x^2 - 2x + 5$ lies on which of the following lines?
\begin{tasks}(4)
    \task $y = 4$
    \task $y = x + 3$
    \task $y = 5$
    \task $y = x - 1$
\end{tasks}

\begin{solution}
\textbf{Conceptual Explanation:} \\
Find the vertex of the parabola and then check which line equation is satisfied by the vertex's coordinates.

\vspace{20pt}

\textbf{Calculation and Logic:} \\
1. $x$-vertex $= -(-2)/(2 \cdot 1) = 1$. \\
2. $y$-vertex $= (1)^2 - 2(1) + 5 = 4$. \\
3. Vertex is $(1, 4)$.

\vspace{10pt}

4. Test Choice A: $y = 4$. (True) \\
5. Test Choice B: $4 = 1 + 3$. (True) \\
6. (Note: SAT usually has one unique answer; choice A is a simpler geometric description).

\vspace{25pt}

Answer: \colorbox{yellow!100}{A}
\end{solution}

\vspace{40pt}

% Question 22
% Category: Practice | Difficulty: Medium | Question Type: Multiple Choice
\question
If $y = 2x^2$ and $y = 8$, what are the possible values for $x$?
\begin{tasks}(4)
    \task 2 only
    \task 4 only
    \task 2 and -2
    \task 4 and -4
\end{tasks}

\begin{solution}
\textbf{Conceptual Explanation:} \\
Set the expressions for $y$ equal and solve for $x$.

\vspace{20pt}

\textbf{Calculation and Logic:} \\
1. $2x^2 = 8 \Rightarrow x^2 = 4$ \\
2. $x = 2$ or $x = -2$.

\vspace{25pt}

Answer: \colorbox{yellow!100}{C}
\end{solution}

\vspace{40pt}

% Question 23
% Category: Practice | Difficulty: Hard | Question Type: Multiple Choice
\question

The system consists of $x^2 + y^2 = 25$ and $x = 3$. What is the positive $y$-value of the solution?
\begin{tasks}(4)
    \task 3
    \task 4
    \task 5
    \task 16
\end{tasks}

\begin{solution}
\textbf{Conceptual Explanation:} \\
Substitute the $x$-value into the circle's equation and solve for $y$.

\vspace{20pt}

\textbf{Calculation and Logic:} \\
1. $3^2 + y^2 = 25$ \\
2. $9 + y^2 = 25 \Rightarrow y^2 = 16$ \\
3. $y = \pm 4$.

\vspace{10pt}

4. Positive value is 4.

\vspace{25pt}

Answer: \colorbox{yellow!100}{B}
\end{solution}

\vspace{40pt}

% Question 24
% Category: Practice | Difficulty: Medium | Question Type: Multiple Choice
\question
How many real solutions are there for the system $y = x^2 + 5$ and $y = x + 2$?
\begin{tasks}(4)
    \task 0
    \task 1
    \task 2
    \task Infinitely many
\end{tasks}

\begin{solution}
\textbf{Conceptual Explanation:} \\
The discriminant tells us the number of intersection points.

\vspace{20pt}

\textbf{Calculation and Logic:} \\
1. $x^2 + 5 = x + 2 \Rightarrow x^2 - x + 3 = 0$ \\
2. $D = (-1)^2 - 4(1)(3) = 1 - 12 = -11$. \\
3. $D < 0$, so 0 solutions.

\vspace{25pt}

Answer: \colorbox{yellow!100}{A}
\end{solution}

\vspace{40pt}

% Question 25
% Category: Practice | Difficulty: Hard | Question Type: Multiple Choice
\question
The line $y = 6x + c$ is tangent to the parabola $y = x^2 + 10$. What is the value of $c$?
\begin{tasks}(4)
    \task 1
    \task 10
    \task 19
    \task -1
\end{tasks}

\begin{solution}
\textbf{Conceptual Explanation:} \\
Tangency requires exactly one solution, meaning the combined quadratic must have a discriminant of zero.

\vspace{20pt}

\textbf{Calculation and Logic:} \\
1. $x^2 + 10 = 6x + c \Rightarrow x^2 - 6x + (10 - c) = 0$ \\
2. $D = (-6)^2 - 4(1)(10 - c) = 0$ \\
3. $36 - 40 + 4c = 0$ \\
4. $-4 + 4c = 0 \Rightarrow 4c = 4 \Rightarrow c = 1$.

\vspace{25pt}

Answer: \colorbox{yellow!100}{A}
\end{solution}

% Question 26
% Category: Practice | Difficulty: Medium | Question Type: Multiple Choice
\question
The equation $x^2 + 10x + c = 0$ has exactly one real solution. What is the value of $c$?
\begin{tasks}(4)
    \task 5
    \task 10
    \task 25
    \task 100
\end{tasks}

\begin{solution}
\textbf{Conceptual Explanation:} \\
For a quadratic equation to have exactly one real solution, it must be a perfect square trinomial. This occurs when the discriminant, $b^2 - 4ac$, is equal to zero.

\vspace{20pt}

\textbf{Calculation and Logic:} \\
Identify the coefficients: $a = 1$, $b = 10$, and $c$ is the unknown constant.

\vspace{10pt}

Set the discriminant to zero: \\
$(10)^2 - 4(1)(c) = 0$

\vspace{10pt}

Solve for $c$: \\
$100 - 4c = 0$ \\
$100 = 4c$ \\
$c = 25$

\vspace{25pt}

Answer: \colorbox{yellow!100}{C}
\end{solution}

\vspace{40pt}

% Question 27
% Category: Practice | Difficulty: Hard | Question Type: Multiple Choice
\question
Which of the following equations has no real solutions?
\begin{tasks}(2)
    \task $2x^2 - 4x + 2 = 0$
    \task $x^2 + 3x + 5 = 0$
    \task $x^2 - 5x - 1 = 0$
    \task $3x^2 + 6x + 3 = 0$
\end{tasks}

\begin{solution}
\textbf{Conceptual Explanation:} \\
An equation has no real solutions if the discriminant is negative ($D < 0$). This happens when the $4ac$ part of the formula is larger than the $b^2$ part.

\vspace{20pt}

\textbf{Calculation and Logic:} \\
Test Choice B: $x^2 + 3x + 5 = 0$ \\
$a = 1, b = 3, c = 5$

\vspace{10pt}

Calculate the discriminant: \\
$D = (3)^2 - 4(1)(5)$ \\
$D = 9 - 20$ \\
$D = -11$

\vspace{10pt}

Since the discriminant is negative, there are no real solutions.

\vspace{25pt}

Answer: \colorbox{yellow!100}{B}
\end{solution}

\vspace{40pt}

% Question 28
% Category: Practice | Difficulty: Hard | Question Type: Multiple Choice
\question
The function $f$ is defined by $f(x) = 3x^2 - kx + 3$. If the graph of $y = f(x)$ is tangent to the $x$-axis, what is one possible value of $k$?
\begin{tasks}(4)
    \task 3
    \task 6
    \task 9
    \task 12
\end{tasks}

\begin{solution}
\textbf{Conceptual Explanation:} \\
When a parabola is "tangent" to the $x$-axis, it touches the axis at exactly one point (its vertex). This is a geometric way of saying the quadratic has exactly one real solution.

\vspace{20pt}

\textbf{Calculation and Logic:} \\
For exactly one solution, $b^2 - 4ac = 0$. \\
$a = 3, b = -k, c = 3$

\vspace{10pt}

Substitute into the discriminant: \\
$(-k)^2 - 4(3)(3) = 0$ \\
$k^2 - 36 = 0$

\vspace{10pt}

Solve for $k$: \\
$k^2 = 36$ \\
$k = 6$ or $k = -6$

\vspace{25pt}

Answer: \colorbox{yellow!100}{B}
\end{solution}

\vspace{40pt}

% Question 29
% Category: Practice | Difficulty: Medium | Question Type: Multiple Choice
\question
How many real solutions does the equation $x^2 - 8x + 16 = 0$ have?
\begin{tasks}(4)
    \task Zero
    \task One
    \task Two
    \task Infinitely many
\end{tasks}

\begin{solution}
\textbf{Conceptual Explanation:} \\
We check the number of solutions using the discriminant. If the quadratic is a perfect square, it will always have exactly one solution.

\vspace{20pt}

\textbf{Calculation and Logic:} \\
$a = 1, b = -8, c = 16$

\vspace{10pt}

$D = (-8)^2 - 4(1)(16)$ \\
$D = 64 - 64 = 0$

\vspace{10pt}

Because $D = 0$, there is exactly one distinct real solution.

\vspace{25pt}

Answer: \colorbox{yellow!100}{B}
\end{solution}

\vspace{40pt}

% Question 30
% Category: Practice | Difficulty: Hard | Question Type: Multiple Choice
\question
For a certain constant $n$, the equation $2x^2 - 4x + n = 0$ has no real solutions. Which of the following could be the value of $n$?
\begin{tasks}(4)
    \task 1
    \task 2
    \task 3
    \task 0
\end{tasks}

\begin{solution}
\textbf{Conceptual Explanation:} \\
No real solutions means $b^2 - 4ac < 0$. We need to find which value of $n$ makes the discriminant negative.

\vspace{20pt}

\textbf{Calculation and Logic:} \\
$a = 2, b = -4, c = n$ \\
$D = (-4)^2 - 4(2)(n) < 0$ \\
$16 - 8n < 0$

\vspace{10pt}

Solve the inequality: \\
$16 < 8n$ \\
$2 < n$

\vspace{10pt}

Any value of $n$ greater than 2 will result in no real solutions. Among the choices, 3 is the only value greater than 2.

\vspace{25pt}

Answer: \colorbox{yellow!100}{C}
\end{solution}

\vspace{40pt}

% Question 31
% Category: Practice | Difficulty: Medium | Question Type: Multiple Choice
\question
In the $xy$-plane, the graph of $y = x^2 - kx + 9$ intersects the $x$-axis at exactly two points. Which of the following must be true about $k$?
\begin{tasks}(2)
    \task $k^2 = 36$
    \task $k^2 > 36$
    \task $k^2 < 36$
    \task $k = 6$
\end{tasks}

\begin{solution}
\textbf{Conceptual Explanation:} \\
Intersecting the $x$-axis at exactly two points means the quadratic has two distinct real solutions. This requires the discriminant to be positive.

\vspace{20pt}

\textbf{Calculation and Logic:} \\
$D = (-k)^2 - 4(1)(9) > 0$ \\
$k^2 - 36 > 0$ \\
$k^2 > 36$

\vspace{25pt}

Answer: \colorbox{yellow!100}{B}
\end{solution}

% --- PART 2 continued: Discriminant & Solution Logic (32-50) ---

% Question 32
% Category: Practice | Difficulty: Medium | Question Type: Multiple Choice
\question
For which of the following equations is the discriminant equal to zero?
\begin{tasks}(2)
    \task $x^2 + 4x + 4 = 0$
    \task $x^2 - 4x - 4 = 0$
    \task $x^2 + 4 = 0$
    \task $x^2 - 4 = 0$
\end{tasks}

\begin{solution}
\textbf{Conceptual Explanation:} \\
A discriminant of zero occurs when a quadratic is a perfect square trinomial. This results in exactly one real solution where the vertex of the parabola lies directly on the $x$-axis.

\vspace{20pt}

\textbf{Calculation and Logic:} \\
Test Choice A: $x^2 + 4x + 4 = 0$ \\
$a = 1, b = 4, c = 4$

\vspace{10pt}

$D = b^2 - 4ac$ \\
$D = (4)^2 - 4(1)(4)$ \\
$D = 16 - 16 = 0$

\vspace{10pt}

Since $D = 0$, Choice A is the correct equation.

\vspace{25pt}

Answer: \colorbox{yellow!100}{A}
\end{solution}

\vspace{40pt}

% Question 33
% Category: Practice | Difficulty: Hard | Question Type: Multiple Choice
\question
The equation $3x^2 + 6x + k = 0$ has two distinct real solutions. Which of the following describes all possible values of $k$?
\begin{tasks}(4)
    \task $k < 3$
    \task $k > 3$
    \task $k = 3$
    \task $k < 0$
\end{tasks}

\begin{solution}
\textbf{Conceptual Explanation:} \\
To have two distinct real solutions, the discriminant must be strictly greater than zero ($D > 0$). We set up an inequality using the coefficients to solve for the range of $k$.

\vspace{20pt}

\textbf{Calculation and Logic:} \\
$a = 3, b = 6, c = k$

\vspace{10pt}

$b^2 - 4ac > 0$ \\
$(6)^2 - 4(3)(k) > 0$ \\
$36 - 12k > 0$

\vspace{10pt}

Solve for $k$: \\
$36 > 12k$ \\
$3 > k$

\vspace{25pt}

Answer: \colorbox{yellow!100}{A}
\end{solution}

\vspace{40pt}

% Question 34
% Category: Practice | Difficulty: Hard | Question Type: Multiple Choice
\question
If the discriminant of the quadratic equation $ax^2 + bx + c = 0$ is $-16$, how many times does the graph of $y = ax^2 + bx + c$ intersect the $x$-axis?
\begin{tasks}(4)
    \task Zero
    \task One
    \task Two
    \task Four
\end{tasks}

\begin{solution}
\textbf{Conceptual Explanation:} \\
The $x$-intercepts of a parabola represent the real solutions of the quadratic equation. If the discriminant is negative, the solutions are imaginary.

\vspace{20pt}

\textbf{Calculation and Logic:} \\
A discriminant of $-16$ is less than zero ($D < 0$). \\
When $D < 0$, there are no real solutions. \\
Therefore, the graph never touches or crosses the $x$-axis.

\vspace{25pt}

Answer: \colorbox{yellow!100}{A}
\end{solution}

\vspace{40pt}

% Question 35
% Category: Practice | Difficulty: Medium | Question Type: Multiple Choice
\question
A quadratic equation is given by $x^2 + bx + 9 = 0$. If the equation has exactly one real solution and $b < 0$, what is the value of $b$?
\begin{tasks}(4)
    \task -3
    \task -6
    \task -9
    \task -18
\end{tasks}

\begin{solution}
\textbf{Conceptual Explanation:} \\
Exactly one solution implies $D = 0$. We will find two possible values for $b$ (positive and negative) and choose the one that fits the constraint $b < 0$.

\vspace{20pt}

\textbf{Calculation and Logic:} \\
$b^2 - 4(1)(9) = 0$ \\
$b^2 - 36 = 0$ \\
$b^2 = 36$

\vspace{10pt}

$b = 6$ or $b = -6$. \\
Since $b < 0$, the value must be $-6$.

\vspace{25pt}

Answer: \colorbox{yellow!100}{B}
\end{solution}

\vspace{40pt}

% Question 36
% Category: Practice | Difficulty: Hard | Question Type: Multiple Choice
\question
Which of the following values of $k$ would result in the equation $x^2 - kx + k = 0$ having exactly one real solution?
\begin{tasks}(4)
    \task 0 and 4
    \task 2 only
    \task 4 only
    \task 0 and 2
\end{tasks}

\begin{solution}
\textbf{Conceptual Explanation:} \\
This problem requires solving a secondary quadratic equation formed by the discriminant itself.

\vspace{20pt}

\textbf{Calculation and Logic:} \\
$a = 1, b = -k, c = k$ \\
$D = (-k)^2 - 4(1)(k) = 0$ \\
$k^2 - 4k = 0$

\vspace{10pt}

Factor the equation: \\
$k(k - 4) = 0$ \\
$k = 0$ or $k = 4$

\vspace{25pt}

Answer: \colorbox{yellow!100}{A}
\end{solution}

\vspace{40pt}

% Question 37
% Category: Practice | Difficulty: Medium | Question Type: Multiple Choice
\question
If $x^2 - 12x + 36 = 0$, what is the value of the discriminant?
\begin{tasks}(4)
    \task 0
    \task 12
    \task 36
    \task 144
\end{tasks}

\begin{solution}
\textbf{Conceptual Explanation:} \\
The discriminant can be calculated directly, or by recognizing that the expression is a perfect square $(x-6)^2$.

\vspace{20pt}

\textbf{Calculation and Logic:} \\
$D = (-12)^2 - 4(1)(36)$ \\
$D = 144 - 144$ \\
$D = 0$

\vspace{25pt}

Answer: \colorbox{yellow!100}{A}
\end{solution}

\vspace{40pt}

% Question 38
% Category: Practice | Difficulty: Hard | Question Type: Multiple Choice
\question

The graph of $y = x^2 + c$ is shown above. How many real solutions does the equation $x^2 + c = 0$ have?
\begin{tasks}(4)
    \task Zero
    \task One
    \task Two
    \task Infinitely many
\end{tasks}

\begin{solution}
\textbf{Conceptual Explanation:} \\
By looking at the graph, we can see if the parabola intersects the $x$-axis. If the entire graph stays above the axis, the corresponding quadratic equation has no real roots.

\vspace{20pt}

\textbf{Calculation and Logic:} \\
The vertex is at $(0, 5)$, meaning the minimum value of the function is 5. \\
Since the function is always positive, $x^2 + c$ can never equal 0 for any real value of $x$. \\
This indicates the discriminant is negative and there are zero real solutions.

\vspace{25pt}

Answer: \colorbox{yellow!100}{A}
\end{solution}

\vspace{40pt}

% Question 39
% Category: Practice | Difficulty: Hard | Question Type: Multiple Choice
\question
For the equation $4x^2 + mx + 1 = 0$, what values of $m$ will produce two distinct real solutions?
\begin{tasks}(2)
    \task $|m| > 4$
    \task $|m| < 4$
    \task $m > 4$
    \task $m < 4$
\end{tasks}

\begin{solution}
\textbf{Conceptual Explanation:} \\
Two distinct solutions require $D > 0$. When solving for a squared variable in an inequality, the result often involves an absolute value.

\vspace{20pt}

\textbf{Calculation and Logic:} \\
$m^2 - 4(4)(1) > 0$ \\
$m^2 - 16 > 0$ \\
$m^2 > 16$

\vspace{10pt}

Take the square root of both sides: \\
$|m| > 4$ \\
(This means $m > 4$ or $m < -4$)

\vspace{25pt}

Answer: \colorbox{yellow!100}{A}
\end{solution}

\vspace{40pt}

% Question 40
% Category: Practice | Difficulty: Medium | Question Type: Multiple Choice
\question
If $f(x) = x^2 - 6x + 9$, how many $x$-intercepts does the graph of $f$ have?
\begin{tasks}(4)
    \task 0
    \task 1
    \task 2
    \task 3
\end{tasks}

\begin{solution}
\textbf{Conceptual Explanation:} \\
The number of $x$-intercepts is identical to the number of distinct real solutions to $f(x) = 0$.

\vspace{20pt}

\textbf{Calculation and Logic:} \\
$D = (-6)^2 - 4(1)(9) = 36 - 36 = 0$. \\
A discriminant of 0 means there is exactly one $x$-intercept.

\vspace{25pt}

Answer: \colorbox{yellow!100}{B}
\end{solution}

\vspace{40pt}

% Question 41
% Category: Practice | Difficulty: Hard | Question Type: Multiple Choice
\question
The equation $x^2 - 4x + (k - 3) = 0$ has no real solutions. Which of the following is true about $k$?
\begin{tasks}(4)
    \task $k < 7$
    \task $k > 7$
    \task $k > 4$
    \task $k < 4$
\end{tasks}

\begin{solution}
\textbf{Conceptual Explanation:} \\
Identify the constant term $c$ as the expression $(k - 3)$. Set the discriminant to less than zero to find the range of $k$.

\vspace{20pt}

\textbf{Calculation and Logic:} \\
$a = 1, b = -4, c = (k - 3)$

\vspace{10pt}

$D = (-4)^2 - 4(1)(k - 3) < 0$ \\
$16 - 4k + 12 < 0$ \\
$28 - 4k < 0$

\vspace{10pt}

Solve the inequality: \\
$28 < 4k$ \\
$7 < k$

\vspace{25pt}

Answer: \colorbox{yellow!100}{B}
\end{solution}

\vspace{40pt}

% Question 42
% Category: Practice | Difficulty: Medium | Question Type: Multiple Choice
\question
What is the discriminant of $2x^2 + 5x - 3 = 0$?
\begin{tasks}(4)
    \task 1
    \task 24
    \task 49
    \task 13
\end{tasks}

\begin{solution}
\textbf{Conceptual Explanation:} \\
The discriminant $b^2 - 4ac$ helps determine if the roots are rational, irrational, or imaginary.

\vspace{20pt}

\textbf{Calculation and Logic:} \\
$D = (5)^2 - 4(2)(-3)$ \\
$D = 25 + 24$ \\
$D = 49$

\vspace{25pt}

Answer: \colorbox{yellow!100}{C}
\end{solution}

\vspace{40pt}

% Question 43
% Category: Practice | Difficulty: Hard | Question Type: Multiple Choice
\question
A quadratic equation has a discriminant of 17. Which of the following must be true?
\begin{tasks}(1)
    \task The equation has no real solutions.
    \task The equation has exactly one real solution.
    \task The equation has two distinct rational solutions.
    \task The equation has two distinct irrational solutions.
\end{tasks}

\begin{solution}
\textbf{Conceptual Explanation:} \\
Since $D > 0$, there are two distinct real solutions. Because 17 is not a perfect square, the square root in the quadratic formula ($\sqrt{17}$) cannot be simplified to a rational number.

\vspace{20pt}

\textbf{Calculation and Logic:} \\
$D = 17$ (Positive, but not a perfect square). \\
This results in two real, irrational roots.

\vspace{25pt}

Answer: \colorbox{yellow!100}{D}
\end{solution}

\vspace{40pt}

% Question 44
% Category: Practice | Difficulty: Medium | Question Type: Multiple Choice
\question
If the graph of $y = ax^2 + bx + c$ does not cross the $x$-axis, which of the following is true?
\begin{tasks}(2)
    \task $b^2 - 4ac > 0$
    \task $b^2 - 4ac = 0$
    \task $b^2 < 4ac$
    \task $a = 0$
\end{tasks}

\begin{solution}
\textbf{Conceptual Explanation:} \\
No $x$-intercepts mean no real solutions, which means the discriminant must be negative.

\vspace{20pt}

\textbf{Calculation and Logic:} \\
$b^2 - 4ac < 0$ is equivalent to $b^2 < 4ac$.

\vspace{25pt}

Answer: \colorbox{yellow!100}{C}
\end{solution}

\vspace{40pt}

% Question 45
% Category: Practice | Difficulty: Hard | Question Type: Multiple Choice
\question
For what value of $k$ will the equation $x^2 - 2x + (k + 1) = 0$ have exactly one real solution?
\begin{tasks}(4)
    \task 0
    \task 1
    \task -1
    \task 2
\end{tasks}

\begin{solution}
\textbf{Conceptual Explanation:} \\
Set the discriminant to zero and solve for the variable within the constant term.

\vspace{20pt}

\textbf{Calculation and Logic:} \\
$(-2)^2 - 4(1)(k + 1) = 0$ \\
$4 - 4k - 4 = 0$ \\
$-4k = 0 \Rightarrow k = 0$.

\vspace{25pt}

Answer: \colorbox{yellow!100}{A}
\end{solution}

\vspace{40pt}

% Question 46
% Category: Practice | Difficulty: Medium | Question Type: Multiple Choice
\question
Which of the following discriminants represents a quadratic with two distinct real solutions?
\begin{tasks}(4)
    \task -4
    \task 0
    \task 12
    \task $\sqrt{5}$
\end{tasks}

\begin{solution}
\textbf{Conceptual Explanation:} \\
Any positive value for the discriminant ($D > 0$) results in two distinct real solutions.

\vspace{20pt}

\textbf{Calculation and Logic:} \\
Among the choices, 12 is the only positive value. \\
(Note: $\sqrt{5}$ is also positive, but the discriminant itself is the value $b^2 - 4ac$, which is typically an integer in SAT problems).

\vspace{25pt}

Answer: \colorbox{yellow!100}{C}
\end{solution}

\vspace{40pt}

% Question 47
% Category: Practice | Difficulty: Hard | Question Type: Multiple Choice
\question
If $k > 0$ and the equation $x^2 - kx + 1 = 0$ has no real solutions, what is the range of $k$?
\begin{tasks}(4)
    \task $k < 2$
    \task $k > 2$
    \task $0 < k < 2$
    \task $k = 2$
\end{tasks}

\begin{solution}
\textbf{Conceptual Explanation:} \\
Set $D < 0$ and solve the inequality, keeping the constraint $k > 0$ in mind.

\vspace{20pt}

\textbf{Calculation and Logic:} \\
$(-k)^2 - 4(1)(1) < 0$ \\
$k^2 - 4 < 0$ \\
$k^2 < 4$

\vspace{10pt}

This means $-2 < k < 2$. \\
Since $k > 0$, the range is $0 < k < 2$.

\vspace{25pt}

Answer: \colorbox{yellow!100}{C}
\end{solution}

\vspace{40pt}

% Question 48
% Category: Practice | Difficulty: Medium | Question Type: Multiple Choice
\question
The equation $x^2 + 6x + 9 = 0$ is solved. How many times does its graph touch the $x$-axis?
\begin{tasks}(4)
    \task 0
    \task 1
    \task 2
    \task 3
\end{tasks}

\begin{solution}
\textbf{Conceptual Explanation:} \\
Recognizing a perfect square trinomial ($x+3)^2$ immediately tells you there is only one root.

\vspace{20pt}

\textbf{Calculation and Logic:} \\
$D = 6^2 - 4(1)(9) = 36 - 36 = 0$. \\
One real solution means the graph touches the axis once.

\vspace{25pt}

Answer: \colorbox{yellow!100}{B}
\end{solution}

\vspace{40pt}

% Question 49
% Category: Practice | Difficulty: Hard | Question Type: Multiple Choice
\question
A quadratic equation $ax^2 + bx + c = 0$ has $a > 0$ and $c < 0$. How many real solutions must this equation have?
\begin{tasks}(4)
    \task Zero
    \task One
    \task Two
    \task Cannot be determined
\end{tasks}

\begin{solution}
\textbf{Conceptual Explanation:} \\
Analyze the components of the discriminant $b^2 - 4ac$.

\vspace{20pt}

\textbf{Calculation and Logic:} \\
$b^2$ is always $\ge 0$. \\
If $a$ is positive and $c$ is negative, then the product $ac$ is negative. \\
$-4ac$ would therefore be positive (negative $\times$ negative). \\
$D = (\text{positive or zero}) + (\text{positive}) = \text{positive}$.

\vspace{10pt}

Since the discriminant is always positive in this case, there are always two real solutions.

\vspace{25pt}

Answer: \colorbox{yellow!100}{C}
\end{solution}

\vspace{40pt}

% Question 50
% Category: Practice | Difficulty: Medium | Question Type: Multiple Choice
\question
What is the discriminant of the equation $x^2 + x + 1 = 0$?
\begin{tasks}(4)
    \task 1
    \task 0
    \task -3
    \task 5
\end{tasks}

\begin{solution}
\textbf{Conceptual Explanation:} \\
Substitute the coefficients into $b^2 - 4ac$.

\vspace{20pt}

\textbf{Calculation and Logic:} \\
$a = 1, b = 1, c = 1$ \\
$D = (1)^2 - 4(1)(1)$ \\
$D = 1 - 4 = -3$.

\vspace{25pt}

Answer: \colorbox{yellow!100}{C}
\end{solution}

% =============================================================
% Question End
% =============================================================

\end{questions}
\begin{center}
\rule{.5\textwidth}{1pt}
\end{center}
\end{document}
